\documentclass[10pt]{extarticle}
\usepackage[a4paper,margin=1.5cm]{geometry}
\usepackage[T1]{fontenc}
\usepackage{amsmath,amssymb}
\usepackage{graphicx}
\usepackage[export]{adjustbox}
\usepackage{multicol}
\usepackage{enumitem}
\usepackage{booktabs,array,tabularx}
\setlength{\columnsep}{0.7cm}
\setlength{\parindent}{0pt}
\setlength{\parskip}{3pt}
\pagestyle{empty}
\makeatletter
\newcommand{\qfigdef}[4]{\@namedef{qf@file@#1}{#2}\@namedef{qf@w@#1}{#3}\@namedef{qf@wide@#1}{#4}}
\newcommand{\qf@check}[1]{\@ifundefined{qf@file@#1}{\PackageError{score2}{Unknown figure id #1}{}}{}}
\newcommand{\qf@img}[4][]{\includegraphics[width=#2\dimexpr\@nameuse{qf@w@#4}pt\relax,max width=#3,max height=0.45\textheight,#1]{\@nameuse{qf@file@#4}}}
\newif\ifqf@full
\newif\ifqf@boxed
\newsavebox{\qf@box}
\newcommand{\qf@open}{\global\qf@fulltrue\par\global\qf@boxedtrue
  \setbox\qf@box\vbox\bgroup\hsize\textwidth\linewidth\textwidth}
\newcommand{\qf@emit}[1]{\par\egroup\global\qf@boxedfalse
  \ifdim\dimexpr\ht\qf@box+\dp\qf@box\relax<#1\textheight\noindent\usebox{\qf@box}\par\else\unvbox\qf@box\fi}
\newcommand{\qfullsplit}{\ifqf@boxed\qf@emit{0.85}\fi}
\newcommand{\qf@close}{\ifqf@boxed\qf@emit{0.45}\fi\global\qf@fullfalse}
\newcommand{\qfullbegin}{\end{multicols}\qf@open}
\newcommand{\qfullend}{\qf@close\begin{multicols}{2}}
\newcommand{\qfullstart}{\qf@open}
\newcommand{\qfulldone}{\qf@close}
\newcommand{\qfig}[2][1]{\qf@check{#2}%
  \ifqf@full\qf@blk{#1}{#2}\else
  \ifnum\@nameuse{qf@wide@#2}=1 \expandafter\@firstoftwo\else\expandafter\@secondoftwo\fi
  {\end{multicols}\begin{center}\qf@img{#1}{\textwidth}{#2}\end{center}\begin{multicols}{2}}%
  {\qf@blk{#1}{#2}}\fi}
\newcommand{\qf@blk}[2]{\par\smallskip{\centering\qf@img{#1}{\linewidth}{#2}\par}\smallskip}
\newcommand{\qfigi}[2][1]{\qf@check{#2}\qf@img[valign=c]{#1}{\linewidth}{#2}}
\makeatother
\newcommand{\papersec}[1]{\par\bigskip{\centering\Large\bfseries PAPER-#1\par}\medskip}
\newcommand{\partsec}[2]{\par\medskip\noindent{\large\bfseries PART-#1: #2}\par\smallskip}
\newcommand{\sect}[1]{\par\smallskip\noindent\textbf{#1}\par}
\newcommand{\qn}[1]{\par\smallskip\noindent\textbf{#1.}~}
\newcommand{\paraq}[2]{\par\smallskip\noindent\textbf{Paragraph for Questions #1 and #2}\par}
\newenvironment{qoptions}{\begin{itemize}[leftmargin=1.8em,itemsep=1pt,topsep=2pt,label={}]}{\end{itemize}}
\newcommand{\opt}[2]{\item[(#1)] #2}
% figure registry (generated by score2_tools.py finalize)
\qfigdef{p1_phy_s1_q01_1}{SRG (Enthusiast) - Test 15 (19 Dec 2021)-img/p1_phy_s1_q01_1.png}{67.6}{0}
\qfigdef{p1_phy_s1_q02_1}{SRG (Enthusiast) - Test 15 (19 Dec 2021)-img/p1_phy_s1_q02_1.png}{108.2}{0}
\qfigdef{p1_phy_s1_q03_1}{SRG (Enthusiast) - Test 15 (19 Dec 2021)-img/p1_phy_s1_q03_1.png}{140.2}{0}
\qfigdef{p1_phy_s1_q04_1}{SRG (Enthusiast) - Test 15 (19 Dec 2021)-img/p1_phy_s1_q04_1.png}{196.4}{0}
\qfigdef{p1_phy_s1_q05_1}{SRG (Enthusiast) - Test 15 (19 Dec 2021)-img/p1_phy_s1_q05_1.png}{119.4}{0}
\qfigdef{p1_phy_s1_q06_1}{SRG (Enthusiast) - Test 15 (19 Dec 2021)-img/p1_phy_s1_q06_1.png}{160.4}{0}
\qfigdef{p1_phy_s2_q07_1}{SRG (Enthusiast) - Test 15 (19 Dec 2021)-img/p1_phy_s2_q07_1.png}{101.6}{0}
\qfigdef{p1_phy_s2_q08_1}{SRG (Enthusiast) - Test 15 (19 Dec 2021)-img/p1_phy_s2_q08_1.png}{178.5}{0}
\qfigdef{p1_phy_s2_q11_1}{SRG (Enthusiast) - Test 15 (19 Dec 2021)-img/p1_phy_s2_q11_1.png}{162.4}{0}
\qfigdef{p1_phy_s2_q12_1}{SRG (Enthusiast) - Test 15 (19 Dec 2021)-img/p1_phy_s2_q12_1.png}{327.3}{1}
\qfigdef{p1_phy_s3_q13_1}{SRG (Enthusiast) - Test 15 (19 Dec 2021)-img/p1_phy_s3_q13_1.png}{224.3}{0}
\qfigdef{p1_phy_s3_q14_1}{SRG (Enthusiast) - Test 15 (19 Dec 2021)-img/p1_phy_s3_q14_1.png}{156.2}{0}
\qfigdef{p1_phy_s3_q15_1}{SRG (Enthusiast) - Test 15 (19 Dec 2021)-img/p1_phy_s3_q15_1.png}{92.1}{0}
\qfigdef{p1_phy_s3_q16_1}{SRG (Enthusiast) - Test 15 (19 Dec 2021)-img/p1_phy_s3_q16_1.png}{118.3}{0}
\qfigdef{p1_phy_s3_q18_1}{SRG (Enthusiast) - Test 15 (19 Dec 2021)-img/p1_phy_s3_q18_1.png}{86.0}{0}

\begin{document}
\begin{center}{\Large\bfseries SRG (Enthusiast) --- Test 15 (19 Dec 2021)}\end{center}
\begin{multicols}{2}
\partsec{1}{PHYSICS}
\sect{Section-I}
\qn{1} In the arrangement shown in figure neglect the masses of the pulley and string, and also friction. The accelerations of blocks A and B are
\qfig{p1_phy_s1_q01_1}
\begin{qoptions}
\opt{A}{$g,\ g/2$}
\opt{B}{$g/2,\ g$}
\opt{C}{$3g/2,\ 3g/4$}
\opt{D}{$g,\ g$}
\end{qoptions}
\qn{2} Given $m_A=30$ kg, $m_B=10$ kg, $m_C=20$ kg. The coefficient of friction between A and B, $\mu_1=0.3$, between B and C, $\mu_2=0.2$ and between C and ground, $\mu_3=0.1$. The least horizontal force $F$ to start motion of any part of the system of three blocks resting upon one another as shown in figure is ($g=10$ m/s$^2$)
\qfig{p1_phy_s1_q02_1}
\begin{qoptions}
\opt{A}{60 N}
\opt{B}{90 N}
\opt{C}{80 N}
\opt{D}{150 N}
\end{qoptions}
\qn{3} Two blocks of masses $m_1$ and $m_2$ are connected with a massless spring and placed over a plank moving with an acceleration `$a$' as shown in figure. The coefficient of friction between the blocks and platform is $\mu$.
\qfig{p1_phy_s1_q03_1}
\begin{qoptions}
\opt{A}{Spring will be stretched if $a>\mu g$}
\opt{B}{Spring will be compressed if $a\le\mu g$}
\opt{C}{Spring will neither be compressed nor be stretched for $a\le\mu g$}
\opt{D}{Spring will be in its natural length under all conditions}
\end{qoptions}
\qn{4} A particle is projected horizontally from a height $b$. The range of values of horizontal speed $v$ so that it falls in the hole of length $a$ is
\qfig{p1_phy_s1_q04_1}
\begin{qoptions}
\opt{A}{$\left(a\sqrt{\dfrac{g}{b}},\,2a\sqrt{\dfrac{g}{b}}\right)$}
\opt{B}{$\left(a\sqrt{\dfrac{g}{2b}},\,a\sqrt{\dfrac{2g}{b}}\right)$}
\opt{C}{$\left(a\sqrt{\dfrac{g}{b}},\,3a\sqrt{\dfrac{g}{b}}\right)$}
\opt{D}{$\left(b\sqrt{\dfrac{g}{b}},\,a\sqrt{\dfrac{g}{b}}\right)$}
\end{qoptions}
\qn{5} A block is in the form of a flat and inclined plane. A body travels with a speed $v_0$ as shown. The minimum value of angle $\beta$ so that the particle does not hit the incline is
\qfig{p1_phy_s1_q05_1}
\begin{qoptions}
\opt{A}{$\dfrac{\pi}{4}$}
\opt{B}{$\dfrac{\pi}{3}$}
\opt{C}{$\tan^{-1}\dfrac{v_0}{\sqrt{2gH}}$}
\opt{D}{$\tan^{-1}\!\left(\dfrac{1}{v_0}\sqrt{\dfrac{gH}{2}}\right)$}
\end{qoptions}
\qn{6} Three small bodies of identical masses can move along a straight line. The central body (2) is initially at rest and bodies 1 and 3 are at a distance $l$ and $2l$ from the central body respectively. Bodies 1 and 3 move towards body 2 with speeds $V_0$ each. The collision between 1 and 2 is perfectly elastic and the collision between body 2 and 3 is perfectly inelastic. After all the collisions are over.
\qfig{p1_phy_s1_q06_1}
\begin{qoptions}
\opt{A}{All the bodies come to rest.}
\opt{B}{The body 1 moves towards left, bodies 2 and 3 move towards right.}
\opt{C}{Body 2 remains at rest and other bodies 1 and 3 turn back.}
\opt{D}{All the bodies move towards right.}
\end{qoptions}
\sect{Section-II}
\qn{7} In the pulley system shown in figure the movable pulleys A, B and C are of mass 1 kg each. D and E are fixed pulleys. The strings are light and inextensible. Choose the correct alternative(s). All pulleys are frictionless.
\qfig{p1_phy_s2_q07_1}
\begin{qoptions}
\opt{A}{Tension in the string is 6.5 N}
\opt{B}{Acceleration of pulley A is $g/3$ downward}
\opt{C}{Acceleration of pulley B is $g/6$ upward}
\opt{D}{Acceleration of pulley C is $g/3$ upward}
\end{qoptions}
\qn{8} The Fig. shows a string of equally placed beads of mass $m$, separated by distance $d$. The beads are free to slide without friction on a thin wire. A constant force $F$ acts on the first bead initially at rest till it makes collision with the second bead. The second bead then collides with the third and so on. Suppose that all collisions are elastic,
\qfig{p1_phy_s2_q08_1}
\begin{qoptions}
\opt{A}{Speed of the first bead immediately before and immediately after its collision with the second bead is $\sqrt{\dfrac{2Fd}{m}}$ and zero respectively.}
\opt{B}{Speed of the first bead immediately before and immediately after its collision with the second bead is $\sqrt{\dfrac{2Fd}{m}}$ and $\dfrac12\sqrt{\dfrac{2Fd}{m}}$ respectively.}
\opt{C}{Speed of the second bead immediately after its collision with third bead is zero.}
\opt{D}{The average speed of the first bead is $\dfrac12\sqrt{\dfrac{2Fd}{m}}$.}
\end{qoptions}
\qn{9} A set of $n$ identical cubical blocks lies at rest parallel to each other along a line on a smooth horizontal surface. The separation between the near surfaces of any two adjacent blocks is $L$. The block at one end is given a speed $v$ towards the next one at time $t=0$. All collisions are completely inelastic, then
\begin{qoptions}
\opt{A}{The last block starts moving at $t=n(n-1)\dfrac{L}{2v}$}
\opt{B}{The last block starts moving at $t=(n-1)\dfrac{L}{v}$}
\opt{C}{The centre of mass of the system will have a final speed $\dfrac{v}{n}$}
\opt{D}{The centre of mass of the system will have a final speed $v$.}
\end{qoptions}
\qn{10} In a one dimensional collision between two identical particles A and B, B is stationary and A has momentum $p$ before impact. During impact, B gives impulse $J$ to A.
\begin{qoptions}
\opt{A}{The total momentum of the `A plus B' system is $p$ before and after the impact, and $(p-J)$ during the impact}
\opt{B}{During the impact A gives impulse $J$ to B}
\opt{C}{The coefficient of restitution is $\dfrac{2J}{p}-1$}
\opt{D}{The coefficient of restitution is $\dfrac{J}{p}+1$}
\end{qoptions}
\qn{11} Two blocks A and B each of mass $m$, are connected by a massless spring of natural length $L$ and spring constant $k$. The blocks are initially resting on a smooth horizontal floor with the spring at its natural length, as shown in fig. A third identical block C, also of mass $m$, moves on the floor with a speed $v$ along the line joining A and B, and collides elastically with A. Then :--
\qfig{p1_phy_s2_q11_1}
\begin{qoptions}
\opt{A}{The kinetic energy of the A--B system, at maximum compression of the spring, is zero}
\opt{B}{The kinetic energy of the A--B system, at maximum compression of the spring, is $\dfrac{mv^2}{4}$}
\opt{C}{The maximum compression of the spring is $v\sqrt{\dfrac{m}{k}}$}
\opt{D}{The maximum compression of the spring is $v\sqrt{\dfrac{m}{2k}}$}
\end{qoptions}
\qfullbegin
\qn{12} Assuming potential energy `$U$' at ground level to be zero.
\qfig{p1_phy_s2_q12_1}
\qfullsplit
All objects are made up of same material.\\
$U_P$ = Potential energy of solid sphere\\
$U_Q$ = Potential energy of solid cube\\
$U_R$ = Potential energy of solid cone\\
$U_S$ = Potential energy of solid cylinder
\begin{qoptions}
\opt{A}{$U_S>U_P$}
\opt{B}{$U_Q>U_S$}
\opt{C}{$U_P>U_Q$}
\opt{D}{$U_S>U_R$}
\end{qoptions}
\qfullend
\sect{Section-III}
\qn{13} Calculate the force of friction (in Newton) for shown situation.
\qfig{p1_phy_s3_q13_1}
\qn{14} Masses $M_1$, $M_2$ and $M_3$ are connected by strings of negligible mass which pass over massless and frictionless pulleys $P_1$ and $P_2$ as shown in fig. The masses move such that the portion of the string between $P_1$ and $P_2$ is parallel to the inclined plane and the portion of the string between $P_2$ and $M_3$ is horizontal. The masses $M_2$ and $M_3$ are 4.0 kg each and the coefficient of kinetic friction between the masses and the surfaces is 0.25. The inclined plane makes an angle of $37^\circ$ with the horizontal. If the mass $M_1$ moves downwards with a uniform velocity, find the mass of $M_1$. (in kg)
\qfig{p1_phy_s3_q14_1}
\qn{15} A string, with one end fixed on a rigid wall, passing over a fixed frictionless pulley at a distance of 2 m from the wall, has a point mass $M=2$ kg attached to it at a distance of 1 m from the wall. A mass $m=0.5$ kg attached at the free end is held at rest so that the string is horizontal between the wall and the pulley and vertical beyond the pulley. What will be the speed (in m/s) with which the mass $M$ will hit the wall when the mass $m$ is released ?
\qfig{p1_phy_s3_q15_1}
\qn{16} A 0.5 kg block slides from the point A (see figure) on a horizontal track with an initial speed of 3 m/s towards a weightless horizontal spring of length 1 m and force constant 2 N/m. The part AB of the track is frictionless and the part BC has the coefficients of static and kinetic friction as 0.22 and 0.2 respectively. If the distance AB and BD are 2 m and 2.14 m respectively, find the total distance (in meter) through which the block moves before it comes to rest completely. (Take $g=10$ m/s$^2$).
\qfig{p1_phy_s3_q16_1}
\qn{17} A ball of mass 100 g is projected vertically upwards from the ground with a velocity of 49 m/s. At the same time another identical ball is dropped from a height of 98 m to fall freely along the same path as that followed by the first ball. After some time the two balls collide and stick together and finally fall to the ground. Find the time of flight (in second) of the masses.
\qn{18} A massless platform is kept on a light elastic spring, shown in Fig. When a sand particle of 0.1 kg mass is dropped on the pan from a height 0.24 m, the particle strikes the pan, and the spring compresses by 0.01 m. From what height (in meter) should the particle be dropped to cause a compression of 0.04 m ?
\qfig{p1_phy_s3_q18_1}
\partsec{2}{CHEMISTRY}
\sect{Section-I}
\qn{1} A mixture containing 0.05 mole of $\mathrm{K_2Cr_2O_7}$ and 0.02 mole of $\mathrm{KMnO_4}$ was treated with excess of KI in acidic medium. The liberated iodine required 1.0 L of $\mathrm{Na_2S_2O_3}$ solution for titration. Concentration of $\mathrm{Na_2S_2O_3}$ solution was :
\begin{qoptions}
\opt{A}{0.40 mol L$^{-1}$}
\opt{B}{0.20 mol L$^{-1}$}
\opt{C}{0.25 mol L$^{-1}$}
\opt{D}{0.30 mol L$^{-1}$}
\end{qoptions}
\qn{2} The energy of separation of an electron in a hydrogen like atom in excited state is 3.4 eV. The de-Broglie wave length (in \AA) associated with the electron is :\\
(Given radius of first orbit of H-atom is 0.53\AA)
\begin{qoptions}
\opt{A}{3.33}
\opt{B}{6.66}
\opt{C}{13.31}
\opt{D}{None of these}
\end{qoptions}
\qn{3} For a real gas (mol mass $=60$) if density at critical point is $0.80\ \mathrm{g/cm^3}$ and its $T_c=\dfrac{4\times10^5}{821}$ K, then van der Waal's constant $a$ (in atm $\mathrm{L^2mol^{-2}}$) is
\begin{qoptions}
\opt{A}{0.3375}
\opt{B}{3.375}
\opt{C}{1.68}
\opt{D}{0.025}
\end{qoptions}
\qn{4} What volume of 75\% alcohol by weight ($d=0.80\,\mathrm{g/cm^3}$) must be used to prepare $150\,\mathrm{cm^3}$ of 30\% alcohol by mass ($d=0.90\,\mathrm{g/cm^3}$) ?
\begin{qoptions}
\opt{A}{67.5 mL}
\opt{B}{56.25 mL}
\opt{C}{44.44 mL}
\opt{D}{32.2 mL}
\end{qoptions}
\qn{5} Using MO theory predict which of the following species has the shortest bond length ?
\begin{qoptions}
\opt{A}{$\mathrm{O_2^{2+}}$}
\opt{B}{$\mathrm{O_2^{+}}$}
\opt{C}{$\mathrm{O_2^{-}}$}
\opt{D}{$\mathrm{O_2^{2-}}$}
\end{qoptions}
\qn{6} The bond dissociation energy of B--F in $\mathrm{BF_3}$ is $646\,\mathrm{kJ\,mol^{-1}}$ where as that of C--F in $\mathrm{CF_4}$ is $515\ \mathrm{kJ\,mol^{-1}}$. The correct reason for higher B--F bond dissociation energy as compared to that of C--F is :
\begin{qoptions}
\opt{A}{Smaller size of B atom as compared to that of C atom}
\opt{B}{Stronger $\sigma$ bond between B and F in $\mathrm{BF_3}$ as compared to that between C and F in $\mathrm{CF_4}$}
\opt{C}{Significant $p\pi-p\pi$ interaction between B and F in $\mathrm{BF_3}$ whereas there is no possibility of such interaction between C and F in $\mathrm{CF_4}$}
\opt{D}{Lower degree of $p\pi-p\pi$ interaction between B and F in $\mathrm{BF_3}$ than that between C and F in $\mathrm{CF_4}$}
\end{qoptions}
\sect{Section-II}
\qn{7} Which of the following statement is / are correct ?
\begin{qoptions}
\opt{A}{$\mathrm{Tl^{3+}}$ salts are oxidising agents}
\opt{B}{$\mathrm{Ga^{+}}$ salts are reducing agents}
\opt{C}{$\mathrm{Pb^{4+}}$ salts are better oxidising agents}
\opt{D}{$\mathrm{As^{5+}}$ salts are oxidising agents}
\end{qoptions}
\qn{8} Which of the following oxides are amphoteric ?
\begin{qoptions}
\opt{A}{BeO}
\opt{B}{SnO}
\opt{C}{ZnO}
\opt{D}{$\mathrm{Al_2O_3}$}
\end{qoptions}
\qn{9} Molecules with see-saw shape are :
\begin{qoptions}
\opt{A}{$\mathrm{XeOF_4}$}
\opt{B}{$\mathrm{SF_4}$}
\opt{C}{$\mathrm{XeOF_2}$}
\opt{D}{$\mathrm{XeO_2F_2}$}
\end{qoptions}
\qn{10} In which of the following species bond angle is expected to be more than $120^\circ$ ?
\begin{qoptions}
\opt{A}{$\mathrm{PCl_4^{+}}$}
\opt{B}{$\mathrm{NO_2^{-}}$}
\opt{C}{$\mathrm{NO_2^{+}}$}
\opt{D}{$\mathrm{XeF_2}$}
\end{qoptions}
\qn{11} Identify correct order of bond angles
\begin{qoptions}
\opt{A}{$\mathrm{Cl_2O>F_2O}$ and $\mathrm{F_2O<H_2O}$}
\opt{B}{$\mathrm{AsI_3>AsBr_3>AsCl_3}$}
\opt{C}{$\mathrm{NO_2^{+}>NO_2^{-}}$}
\opt{D}{$\angle \mathrm{H_bBH_b}>\angle \mathrm{H_tBH_t}$; where $\mathrm{H_t}$ is terminal Hydrogen of $\mathrm{B_2H_6}$ and $\mathrm{H_b}$ is the bridging Hydrogen of $\mathrm{B_2H_6}$}
\end{qoptions}
\qn{12} Which among the following compounds are paramagnetic ?
\begin{qoptions}
\opt{A}{$\mathrm{KO_2}$}
\opt{B}{$\mathrm{K_2O_2}$}
\opt{C}{$\mathrm{K_2O}$}
\opt{D}{$\mathrm{NO_2}$}
\end{qoptions}
\sect{Section-III}
\qn{13} 2 mole, equimolar mixture of $\mathrm{Na_2C_2O_4}$ and $\mathrm{H_2C_2O_4}$ required $V_1$ L of 0.1 M $\mathrm{KMnO_4}$ in acidic medium for complete oxidation. The same amount of the mixture required $V_2$ L of 0.2 M NaOH for neutralization. Then the value of $(V_1+V_2)$ is
\qn{14} If I excitation energy for the H-like (hypothetical) sample is 24 eV, then binding energy in III excited state in eV is :
\qn{15} The van der Waal's constant `b' of a gas is $4\pi\times10^{-4}$ L/mol the centres of the two molecules approach each other is $10^{-x}$ m. What is $x$ ? $N_A=6\times10^{23}$
\qn{16} The wave function of atomic orbital of H like atoms is given as under
\[\psi_{2s}=\frac{1}{4\sqrt{2\pi}}\,z^{3/2}\,[\,2-Zr\,]\,e^{Zr/2}\]
Given that the radius is in \AA, then radius of nodal surface for $\mathrm{He^+}$ ion is (in \AA)
\qn{17} The first ionisation energy of Li is 5.41 eV and electron affinity of chlorine is 3.61 eV. The $\Delta H$ in $\mathrm{kJ\,mol^{-1}}$ for the reaction :\\
$\mathrm{Li(g)+Cl(g)\to Li^{+}(g)+Cl^{-}(g)}$ is $1.737\times10^{x}\ \mathrm{kJ\,mol^{-1}}$\\
What is the value of $x$ ?
\qn{18} Among the triatomic molecules / ions, $\mathrm{BeCl_2,\ N_3^{-},\ N_2O,\ NO_2^{+},\ O_3,\ SCl_2,\ ICl_2^{-},\ I_3^{-}}$ and $\mathrm{XeF_2}$, the total number of linear molecule(s) / ion(s) where the hybridization of the central atom does not have contribution from the d-orbital(s) is :
\partsec{3}{MATHEMATICS}
\sect{Section-I}
\qn{1} If terms independent of $x$ in the expansion of $\left(3x-\dfrac1x\right)^{20}$ and $\left(x+\dfrac{\sqrt[9]{3^{10}}}{x}\right)^{18}$ are $A$ and $B$ respectively then $\left(\dfrac{9}{38}A+B\right)$ equals
\begin{qoptions}
\opt{A}{$3^{10}\cdot{}^{19}C_8$}
\opt{B}{$3^{10}\cdot{}^{19}C_9$}
\opt{C}{$3^{9}\cdot{}^{20}C_8$}
\opt{D}{$3^{9}\cdot{}^{19}C_{10}$}
\end{qoptions}
\qn{2} There are 3 women participating in a chess tournament. Every participant played 3 games with other participants. Number of games that the men played among themselves proved to exceed by 27 number of games that men played with women. Then number of men playing in the tournament is \underline{\hspace{1cm}}
\begin{qoptions}
\opt{A}{3}
\opt{B}{5}
\opt{C}{9}
\opt{D}{12}
\end{qoptions}
\qn{3} The least integer `$a$' for which $1+\log_3\left(x^2+1\right)\le\log_3\left(ax^2+4x+a\right)$ is true for all real values of $x$ is
\begin{qoptions}
\opt{A}{3}
\opt{B}{4}
\opt{C}{5}
\opt{D}{6}
\end{qoptions}
\qn{4} The least value of $x$ satisfying $\dfrac{1}{[x]}+\dfrac{1}{[2x]}=\{x\}+\dfrac13$ is $N$, then number of positive divisors of $24N$ is\\
(where $[\cdot]$ is the greatest integer function and $\{x\}$ is fractional part of $x$)
\begin{qoptions}
\opt{A}{2}
\opt{B}{4}
\opt{C}{6}
\opt{D}{8}
\end{qoptions}
\qn{5} If $S_n(\theta)=\displaystyle\sum_{r=1}^{n}\dfrac{\cos(r-1)\theta}{(\cos\theta)^{r-1}}$, then sum of the series
\[\left[S_3\!\left(\tfrac{\pi}{3}\right)\right]+\left[S_9\!\left(\tfrac{\pi}{3}\right)\right]+\left[S_{15}\!\left(\tfrac{\pi}{3}\right)\right]+\dots+\left[S_{2025}\!\left(\tfrac{\pi}{3}\right)\right]\]
is\\
(where $[\cdot]$ denotes the greatest integer function)
\begin{qoptions}
\opt{A}{0}
\opt{B}{1}
\opt{C}{3}
\opt{D}{2025}
\end{qoptions}
\qn{6} If $2\cos^4x+\sqrt5\left(\sqrt5-1\right)=\cos^2x\left(\sqrt5+1+2\sin x\right)+\left(\sqrt5-1\right)\sin x$, then $\sin^{10}x+4\sin^9x+6\sin^8x+4\sin^7x+\sin^6x+\sin x+1$ equals
\begin{qoptions}
\opt{A}{1}
\opt{B}{2}
\opt{C}{3}
\opt{D}{4}
\end{qoptions}
\sect{Section-II}
\qn{7} In a game a person P at $(x,y)$ can jump to any one of the following points relatively $(x+2,y)$, $(x+4,y)$, $(x,y+2)$, $(x,y+4)$. If P is at $(1,1)$ then the number of ways to reach the point.
\begin{qoptions}
\opt{A}{$(5,5)$ is 14}
\opt{B}{$(5,5)$ is 24}
\opt{C}{$(7,8)$ is 0}
\opt{D}{$(9,9)$ through $(5,5)$ is 196}
\end{qoptions}
\qn{8} If the equation $\left(2+\sin^2x\right)^2-(a-3)\left(2+\sin^2x\right)\left(1+\sin^2x\right)+(a-4)\left(1+\sin^2x\right)^2=0$ has atleast one real root then the complete set of values of $a$ is $[\alpha,\beta]$ then
\begin{qoptions}
\opt{A}{$\alpha+\beta=11$}
\opt{B}{$\beta-\alpha=\dfrac12$}
\opt{C}{$\alpha+\beta=\dfrac{33}{2}$}
\opt{D}{$\beta-\alpha=1$}
\end{qoptions}
\qn{9} Let $S=\displaystyle\sum_{k=1001}^{3001}\left(\dfrac1k\right)$, then
\begin{qoptions}
\opt{A}{$S>1$}
\opt{B}{$S<1$}
\opt{C}{$S>\dfrac43$}
\opt{D}{$S<\dfrac43$}
\end{qoptions}
\qn{10} Let $N=2^2\times3^3\times7^4$, then
\begin{qoptions}
\opt{A}{Number of factors of $N^2$ which are less than $N$ and are not the factors of $N$ is 98}
\opt{B}{Number of factors of $N$ which are of the form $4K+1$, $K\in\mathbb N$ is 9}
\opt{C}{Number of factors of $N$ which are multiples of 54 is 10}
\opt{D}{Number of factors of $N$ which are of the form $4K+1$, $K\in\mathbb N$ is 8}
\end{qoptions}
\qn{11} The series $\dfrac85+\dfrac{16}{65}+\dfrac{24}{325}+\dots\dots\dfrac{8n}{4n^4+1}$
\begin{qoptions}
\opt{A}{Sum of infinite number of terms of series is 4}
\opt{B}{Sum of $n$ terms of series is less than 2 $\forall\ n\in\mathbb N$}
\opt{C}{Sum to $n$ terms of series cannot be an integer for any $n\in\mathbb N$}
\opt{D}{Sum of infinite number of terms of series is less than 3}
\end{qoptions}
\qn{12} Consider 5 points in plane are situated so that no two of the straight lines joining them are parallel, perpendicular, or co incident. From each point perpendiculars the drawn to all the lines joining the other four points. If $k$ is the maximum number of intersections that these perpendiculars can have ? Then
\begin{qoptions}
\opt{A}{More than 310}
\opt{B}{$k$ is divisible by 5}
\opt{C}{$k$ is less than 400}
\opt{D}{$k$ is even number}
\end{qoptions}
\sect{Section-III}
\qn{13} If $S=\displaystyle\lim_{n\to\infty}\sum_{r=2}^{n}\dfrac{{}^rC_2}{3^r}$. Then value of $48S$ is
\qn{14} If $\alpha,\ \beta$ are two roots of the equation $1!+2!+3!+\dots\dots(x-1)!+x!=k^2$, $k\in I^+$. Where $\alpha<\beta$ and if $\alpha_1,\ \alpha_2,\ \alpha_3,\ \alpha_4$ are roots of equation $\left(a+\sqrt b\right)^{x^2-\left(\alpha+3a+5\alpha^2\right)}+\left(a-\sqrt b\right)^{x^2-3\beta}=2a$, where $a^2-b=1$. Then the value of $|\alpha_1+\alpha_2+\alpha_2+\alpha_3+\alpha_1\alpha_2\alpha_3\alpha_4|$ is equal to
\qn{15} The real value of $x$ satisfying $\sqrt[3]{20x+\sqrt[3]{20x+13}}=13$ can be expressed as $\dfrac ab$ where $a$ and $b$ are relatively prime positive integers. Then the value of $(a-100b)$ is
\qn{16} Number of order pairs $(x,y)$ which satisfy the equation $x^2+2x\sin(xy)+1=0$, where $y\in[0,2\pi]$ is
\qn{17} Let $P(x)=x^6-x^5-x^3-x^2-x$ and $\alpha,\ \beta,\ \gamma,\ \delta$ are the roots of the equation $x^4-x^3-x^2-1=0$, then $P(\alpha)+P(\beta)+P(\gamma)+P(\delta)=$
\qn{18} The result of 22 hockey matches are to be predicted. Given that result can be win, lose or draw. If $p$ is number of different forecast which contain exactly 18 correct results and $q$ is number of different forecast which contains exactly 19 correct results. Then value of $\left(\dfrac{4p}{q}\right)$ is
\end{multicols}
\end{document}
